%%%%%%%%%%%%%%%%%%%%%%%%%
\section{\label{sec:configuration}Description of a configuration with particles, coordinates and periodic boundaries}
%%%%%%%%%%%%%%%%%%%%%%%%%

In this article, a configuration refers to the particles, coordinates and periodic boundaries.
For an $n$-component system, the total number of particles, $N=\sum_{i}^{n} N_i$, where $N_i$ is the number of particles of type $i$.
In bulk systems, there are no interactions of the particles with boundaries.
To avoid spurious effects of walls or a container, and thereby minimizing the size of the simulated system needed to model the bulk, periodic boundary conditions are imposed such that particles interact with their nearest neighboring image \cite{allen_computer_1989}.
To simplify the notation in the following derivations, the periodic boundaries align with the Cartesian axes to form a cuboid of $D$ linear dimensions $L_d$, where the volume, $V$, is given by $V=\prod_d^D L_d$.

The position of each particle is represented by a single vector coordinate $\mathbf{r}$, and all internal and/or orientational coordinates of the particle are represented by a single vector $\boldsymbol{\omega}$.
A particle may be composed of multiple interaction sites that are mutually bound, flexibly or rigidly or both, and the site-site interactions could be orientationally-dependent (e.g., point dipoles).
The orientation of a particle might be given implicitly (not part of $\boldsymbol{\omega}$) via the Cartesian coordinates of its interaction sites.
Alternatively, it may be specified explicitly via an appropriate orientation vector, with the location of the interaction sites (if a multi-site particle) constructed on-the-fly from the position, orientation, and internal coordinates.
In a 3-dimensional space, the orientation coordinate for a cylindrically-symmetric object requires just two independent values, e.g., a unit  vector describing the cylindrical axis direction.
For particles with no such symmetry, three independent values are needed to specify the orientation.
Euler angles are sometimes invoked for this, but they are not particularly useful in the context of molecular simulation.
A better choice is a quaternion, which is a normalized, four-element vector \cite{vesely_angular_1982, karney_quaternions_2007}.

If the particle is not isotropic or rigid, the derivations in this article leave a placeholder variable for the orientation contributions that are not further detailed and are beyond the scope of this article.
